% GENERATED FILE. Edit canonical lessons, then run npm run generate:books.
% canonical-source-hash: fnv1a64:2886bd7ec42b0dae
% canonical-lessons: GU-C199-bhid, GU-C199-katar, GU-C199-niyam, GU-C199-ghusvu, GU-C199-khasvu

\chapter{Queues and crowds}
\label{ch:gu-queues-and-crowds}
\hlchaptermodality{199}

\begin{chapteropening}
\textbf{By the end of this chapter:} \emph{I can wait in a queue and say who pushes in.}

Complete the last lesson of chapter 199: I can wait in a queue and say who pushes in.
\end{chapteropening}

\section[\texorpdfstring{\gu{ભીડ} (bhīḍ)}{ભીડ (bhīḍ)}]{\gu{ભીડ} (bhīḍ) — crowd}

\label{lesson:GU-C199-bhid}



\begin{quote}
Before the new one: say the Gujarati for ban, then the Gujarati for danger.
\end{quote}

\subsection*{You'll want to know: \gu{ભીડ}}
\textbf{\gu{ભીડ}} — \emph{bhīḍ} — \textquotedblleft{}crowd\textquotedblright{}.

\begin{quote}
\textbf{\gu{બજારમાં} \gu{ઘણી} \gu{ભીડ} \gu{છે}.} — \emph{bajārmā̃ ghaṇī bhīḍ che} — \textquotedblleft{}The market is very crowded.\textquotedblright{}
\end{quote}

\begin{cousinweb}
\textbf{\gu{ભીડ}} (\emph{bhīḍ}) is a tadbhava. CDIAL lists it, glossed \textquotedblleft{}crowd\textquotedblright{}, under \emph{*bhiṭ}, \textquotedblleft{}meet, throng\textquotedblright{}, a reconstructed form (CDIAL 9490).
\end{cousinweb}

\begin{grammarlens}[title={What you've built}]
One more word for everyday life at A2.
\end{grammarlens}

\subsection*{Your turn}
\begin{itemize}
\raggedright
  \item \emph{Say it:} \emph{bhīḍ}
  \item \emph{Say it:} \emph{bhīḍ}, once more
  \item \emph{Say it:} \emph{bajārmā̃ ghaṇī bhīḍ che} — The market is very crowded.
  \item \emph{Recall:} say the Gujarati for ban, then the Gujarati for danger, then say \emph{bhīḍ} again
\end{itemize}

\begin{tcolorbox}[breakable,colback=teal!4,colframe=teal!35!black,title={Before you move on}]
What does \gu{ભીડ} mean? (\textquotedblleft{}Crowd\textquotedblright{}.) Say it once more.
\end{tcolorbox}
\section[\texorpdfstring{\gu{કતાર} (katār)}{કતાર (katār)}]{\gu{કતાર} (katār) — queue, line}

\label{lesson:GU-C199-katar}



\begin{quote}
Before the new one: say the Gujarati for danger, then the Gujarati for crowd.
\end{quote}

\subsection*{You'll want to know: \gu{કતાર}}
\textbf{\gu{કતાર}} — \emph{katār} — \textquotedblleft{}queue, line\textquotedblright{}.

\begin{quote}
\textbf{\gu{બેંકમાં} \gu{લાંબી} \gu{કતાર} \gu{છે}.} — \emph{bẽkmā̃ lā̃bī katār che} — \textquotedblleft{}There is a long queue at the bank.\textquotedblright{}
\end{quote}

\begin{cousinweb}
\textbf{\gu{કતાર}} (\emph{katār}) is a loan from Arabic \emph{qiṭār}, \textquotedblleft{}a line of camels\textquotedblright{}, which reached Gujarati through Persian. Its Arabic root is \emph{q-ṭ-r} (Lane).
\end{cousinweb}

\begin{grammarlens}[title={What you've built}]
One more word for everyday life at A2.
\end{grammarlens}

\subsection*{Your turn}
\begin{itemize}
\raggedright
  \item \emph{Say it:} \emph{katār}
  \item \emph{Say it:} \emph{katār}, once more
  \item \emph{Say it:} \emph{bẽkmā̃ lā̃bī katār che} — There is a long queue at the bank.
  \item \emph{Recall:} say the Gujarati for danger, then the Gujarati for crowd, then say \emph{katār} again
\end{itemize}

\begin{tcolorbox}[breakable,colback=teal!4,colframe=teal!35!black,title={Before you move on}]
What does \gu{કતાર} mean? (\textquotedblleft{}Queue, line\textquotedblright{}.) Say it once more.
\end{tcolorbox}
\section[\texorpdfstring{\gu{નિયમ} (niyam)}{નિયમ (niyam)}]{\gu{નિયમ} (niyam) — rule, regulation}

\label{lesson:GU-C199-niyam}



\begin{quote}
Before the new one: say the Gujarati for crowd, then the Gujarati for queue.
\end{quote}

\subsection*{You'll want to know: \gu{નિયમ}}
\textbf{\gu{નિયમ}} — \emph{niyam} — \textquotedblleft{}rule, regulation\textquotedblright{}.

\begin{quote}
\textbf{\gu{આ} \gu{બેંકનો} \gu{નિયમ} \gu{છે}.} — \emph{ā bẽkno niyam che} — \textquotedblleft{}This is the bank's rule.\textquotedblright{}
\end{quote}

\begin{cousinweb}
\textbf{\gu{નિયમ}} (\emph{niyam}) is Sanskrit \emph{niyama}, \textquotedblleft{}restraint, rule\textquotedblright{}, taken over whole as a learned loan, a tatsama. CDIAL lists the verb \emph{niyamayati}, \textquotedblleft{}restrains\textquotedblright{} (CDIAL 7251).
\end{cousinweb}

\begin{grammarlens}[title={What you've built}]
One more word for everyday life at A2.
\end{grammarlens}

\subsection*{Your turn}
\begin{itemize}
\raggedright
  \item \emph{Say it:} \emph{niyam}
  \item \emph{Say it:} \emph{niyam}, once more
  \item \emph{Say it:} \emph{ā bẽkno niyam che} — This is the bank's rule.
  \item \emph{Recall:} say the Gujarati for crowd, then the Gujarati for queue, then say \emph{niyam} again
\end{itemize}

\begin{tcolorbox}[breakable,colback=teal!4,colframe=teal!35!black,title={Before you move on}]
What does \gu{નિયમ} mean? (\textquotedblleft{}Rule, regulation\textquotedblright{}.) Say it once more.
\end{tcolorbox}
\section[\texorpdfstring{\gu{ઘૂસવું} (ghūsvũ)}{ઘૂસવું (ghūsvũ)}]{\gu{ઘૂસવું} (ghūsvũ) — to sneak in, to push in}

\label{lesson:GU-C199-ghusvu}



\begin{quote}
Before the new one: say the Gujarati for queue, then the Gujarati for rule.
\end{quote}

\subsection*{You'll want to know: \gu{ઘૂસવું}}
\textbf{\gu{ઘૂસવું}} — \emph{ghūsvũ} — \textquotedblleft{}to sneak in, to push in\textquotedblright{}.

\begin{quote}
\textbf{\gu{એક} \gu{માણસ} \gu{કતારમાં} \gu{ઘૂસે} \gu{છે}.} — \emph{ek māṇas katārmā̃ ghūse che} — \textquotedblleft{}A man pushes into the queue.\textquotedblright{}
\end{quote}

\begin{cousinweb}
\textbf{\gu{ઘૂસવું}} (\emph{ghūsvũ}) is a tadbhava. CDIAL lists it, glossed \textquotedblleft{}to enter forcibly\textquotedblright{}, under \emph{*ghuss}, \textquotedblleft{}thrust in, pierce\textquotedblright{}, a reconstructed form (CDIAL 4492).
\end{cousinweb}

\begin{grammarlens}[title={What you've built}]
One more word for everyday life at A2.
\end{grammarlens}

\subsection*{Your turn}
\begin{itemize}
\raggedright
  \item \emph{Say it:} \emph{ghūsvũ}
  \item \emph{Say it:} \emph{ghūsvũ}, once more
  \item \emph{Say it:} \emph{ek māṇas katārmā̃ ghūse che} — A man pushes into the queue.
  \item \emph{Recall:} say the Gujarati for queue, then the Gujarati for rule, then say \emph{ghūsvũ} again
\end{itemize}

\begin{tcolorbox}[breakable,colback=teal!4,colframe=teal!35!black,title={Before you move on}]
What does \gu{ઘૂસવું} mean? (\textquotedblleft{}To sneak in, to push in\textquotedblright{}.) Say it once more.
\end{tcolorbox}
\section[\texorpdfstring{\gu{ખસવું} (khasvũ)}{ખસવું (khasvũ)}]{\gu{ખસવું} (khasvũ) — to move aside, to shift}

\label{lesson:GU-C199-khasvu}



\begin{quote}
Before the new one: say the Gujarati for rule, then the Gujarati for to sneak in.
\end{quote}

\subsection*{You'll want to know: \gu{ખસવું}}
\textbf{\gu{ખસવું}} — \emph{khasvũ} — \textquotedblleft{}to move aside, to shift\textquotedblright{}.

\begin{quote}
\textbf{\gu{તે} \gu{બાજુમાં} \gu{ખસ્યો}.} — \emph{te bājumā̃ khasyo} — \textquotedblleft{}He moved aside.\textquotedblright{}
\end{quote}

\begin{cousinweb}
\textbf{\gu{ખસવું}} (\emph{khasvũ}) is a tadbhava. CDIAL lists it, glossed \textquotedblleft{}to slide, move, go to a distance\textquotedblright{}, under \emph{*khasati}, \textquotedblleft{}falls, slips\textquotedblright{}, a reconstructed form (CDIAL 3856).
\end{cousinweb}

\begin{grammarlens}[title={What you've built}]
One more word for everyday life at A2.
\end{grammarlens}

\subsection*{Your turn}
\begin{itemize}
\raggedright
  \item \emph{Say it:} \emph{khasvũ}
  \item \emph{Say it:} \emph{khasvũ}, once more
  \item \emph{Say it:} \emph{te bājumā̃ khasyo} — He moved aside.
  \item \emph{Recall:} say the Gujarati for rule, then the Gujarati for to sneak in, then say \emph{khasvũ} again
\end{itemize}

\begin{tcolorbox}[breakable,colback=teal!4,colframe=teal!35!black,title={Before you move on}]
What does \gu{ખસવું} mean? (\textquotedblleft{}To move aside, to shift\textquotedblright{}.) Say it once more.
\end{tcolorbox}
